\color{unidarkgreen}{\subsection{Газовая защита}}
\color{black}

\begin{enumerate}

\item
Газовое реле трансформатора (в составе ГЗ АТ, ГЗ ЛРТ) реагирует на выделение газа при повреждениях внутри кожуха, а также на понижение уровня масла. Для защиты устройства РПН применяется струйное реле (в составе ГЗ РПН), которое реагирует на ускоренный поток масла, возникающий при повреждениях в баке РПН.

\item 
Уставки ГЗ АТ приведены в таблицах \ref{app_set:gzotkl}, \ref{app_set:gzsign}. Уставки ГЗ ЛРТ приведены в таблицах \ref{app_set:gzlrtotkl}, \ref{app_set:gzlrtsign}. Уставки ГЗ РПН приведены в таблице \ref{app_set:gzrpn}.

\item
Предусмотрена возможность приема сигналов срабатывания сигнального контакта ГЗ АТ, отключающего контакта ГЗ АТ и низкой изоляции цепи отключающего контакта ГЗ АТ, по разным цепям от двух источников ПДС1 и ПДС2. Для приема сигналов от ПДС1 предусмотрены входные сигналы \textit{«Ср.сигн контакта ГЗ АТ 1 цепь»}, \textit{«Ср. откл контакта ГЗ АТ 1 цепь»} и \textit{«Низкая изол. ГЗ АТ откл 1 цепь»}, а для приема сигналов от ПДС2 предусмотрены входные сигналы \textit{«Ср. сигн контакта ГЗ АТ 2 цепь»}, \textit{«Ср. откл контакта ГЗ АТ 2 цепь»} и \textit{«Низкая изол. ГЗ АТ откл 2 цепь»}. Если сигналы срабатывания ГЗ АТ заводятся непосредственно со шкафа ШОМ силового трансформатора, то конфигурируются только входные сигналы \textit{«Ср. сигн контакта ГЗ АТ 1 цепь»}, \textit{«Ср. откл контакта ГЗ АТ 1 цепь»} и \textit{«Низкая изол. ГЗ АТ откл 1 цепь»}, а действие ГЗ АТ по второй цепи сигнального и отключающего контакта выводится из работы с помощью уставок \textbf{«ГЗ АТ сигн 2 цепь»} и \textbf{«ГЗ АТ откл 2 цепь»}.

\item
Ввод в работу ГЗ АТ осуществляется с помощью уставок \textbf{«ГЗ АТ сигн 1 цепь»}, \textbf{«ГЗ АТ сигн 2 цепь»}, \textbf{«ГЗ АТ откл 1 цепь»} и \textbf{«ГЗ АТ откл 2 цепь»}. Оперативный перевод ГЗ АТ на сигнал производится с помощью команд управления \textit{«Перевод ГЗ АТ откл на сигнал»}. Предусматривается возможность действия срабатывания сигнального контакта ГЗ АТ на отключение трансформатора с помощью уставки
\textbf{«Откл АТ от ГЗ сигн»}.

\item
Оперативный ввод/вывод сигнальной ступени ГЗ АТ производится с помощью БУ (блока управления функциями) \textit{«Вывод ГЗ АТ сигн»}, а оперативный ввод/вывод отключающей ступени ГЗ АТ  производится с помощью БУ (блока управления функциями) \textit{«Вывод ГЗ АТ откл»}.

\item
Непрерывный контроль изоляции цепи отключающего контакта ГЗ АТ осуществляется с помощью реле контроля изоляции (РКИ). В случае снижения сопротивления изоляции РКИ подает сигнал \textit{«Низкая изол. ГЗ АТ откл 1 цепь»} (\textit{«Низкая изол. ГЗ АТ откл 2 цепь»}) для блокировки действия ГЗ АТ по цепи отключающего контакта. Время срабатывания блокировки от сигналов РКИ задается с помощью уставки \textbf{«Тбл ГЗ АТ откл»}. Ввод блокировки ГЗ АТ от РКИ осуществляется с помощью уставки \textbf{«Блок ГЗ АТ откл»}. Для сброса блокировки от РКИ, используется команда управления \textit{«Возврат ГЗ АТ откл»}.

\item
Алгоритм работы сигнальной ступени ГЗ АТ приведен на рисунке \ref{pic:GZAT_sign}, алгоритм работы отключающей ступени ГЗ АТ приведен на рисунке \ref{pic:GZAT_otkl}. Алгоритмы работы сигнальной и отключающей ступеней ГЗ ЛРТ выполнены аналогично соответствующим ступеням ГЗ АТ.

\begin{figure}[H]
\centering
\includegraphics[width=0.6\textwidth,height=0.6\textheight,keepaspectratio]{\fbpath/06.02 ДЗАТ/041.1071 ГЗ АТ/img/GZAT_otkl.pdf}
\captionof{figure}{Функциональная схема логики отключающей ступени ГЗ АТ}\label{pic:GZAT_otkl}
\end{figure}

\vspace{2mm}
\begin{figure}[H]
\centering
\includegraphics[width=0.65\textwidth,height=0.65\textheight,keepaspectratio]{\fbpath/06.02 ДЗАТ/041.1071 ГЗ АТ/img/GZAT_sign.pdf}
\captionof{figure}{Функциональная схема логики сигнальной ступени ГЗ АТ}\label{pic:GZAT_sign}
\end{figure}

\item
Алгоритм работы ГЗ РПН приведен на рисунке \ref{pic:GZRPN}.
\vspace{2mm}
\begin{figure}[H]
\centering
\includegraphics[width=1\textwidth,height=1\textheight,keepaspectratio]{\fbpath/06.02 ДЗАТ/041.1071 ГЗ АТ/img/GZ_RPN.pdf}
\captionof{figure}{Функциональная схема логики ГЗ РПН}\label{pic:GZRPN}
\end{figure}






\end{enumerate}